% © 2026 Ing. David Jaroš — CC BY-NC-ND 4.0
% UBT-AI-PROVENANCE-BEGIN
% schema: ubt-ai-provenance/v1
% tier: B_machine_verified
% ai_assistance: disclosed
% human_review: machine-verification
% editorial_responsibility: Ing. David Jaroš
% policy: ../../AI_PROVENANCE.md
% notice: Same-carrier commutant claim machine-verified; physical routing implications remain model-level.
% UBT-AI-PROVENANCE-END

\documentclass[11pt,a4paper]{article}
\usepackage[margin=2.5cm]{geometry}
\usepackage{amsmath,amssymb,amsthm}
\usepackage[most]{tcolorbox}
\newtheorem{theorem}{Theorem}
\newtheorem{corollary}[theorem]{Corollary}
\newtheorem{remark}[theorem]{Remark}
\title{Lorentz--Colour Compatibility in UBT:\\
A Same-Carrier Commutant No-Go}
\author{Ing.\ David Jaro\v{s}}
\date{9 October 2026}
\begin{document}
\maketitle

\begin{abstract}
We test whether the algebraic $SU(3)$ colour action can be identified directly
with an internal Standard-Model symmetry acting on the same
$W=\mathbb C\otimes\mathbb H\simeq M_2(\mathbb C)$ carrier that transforms
under the canonical Lorentz action $X\mapsto SXS^\dagger$.
The complex-linear commutant of the full Lorentz action on this carrier is
one-dimensional: only scalar multiples of the identity commute with all six
Lorentz generators.  Therefore a nontrivial $SU(3)$ internal symmetry cannot
act on this same matrix index while commuting with Lorentz transformations.
This does not invalidate the algebraic $SU(3)$ stabiliser theorem; it fixes the
physical bridge requirement.  Colour must act on a distinct internal
multiplicity/mode bundle, or on another action-derived Lorentz-scalar carrier,
rather than being identified naively with the spacetime biquaternion index.
\end{abstract}

\section{Setup}

Let
\[
 W=M_2(\mathbb C)\simeq\mathbb C^4.
\]
The Lorentz action used in the UBT action audit is
\[
 \rho(S)X=SXS^\dagger,\qquad S\in SL(2,\mathbb C).
\]
As a complex representation of the real Lorentz group this is the complexified
four-vector carrier.

Choose the ordered complex basis
\[
 (\mathbf1,I,J,K).
\]
Rotations act on the spatial triplet $(I,J,K)$ and boosts mix the scalar
direction with the corresponding spatial direction.

\section{Commutant theorem}

\begin{theorem}[Same-carrier Lorentz commutant]
Let $C\in\operatorname{End}_{\mathbb C}(W)$ commute with the infinitesimal
action of all three rotations and all three boosts.  Then
\[
 \boxed{C=c\,\mathbf1_W,\qquad c\in\mathbb C.}
\]
Equivalently,
\[
 \operatorname{End}_{\mathrm{Lorentz}}(W)\simeq\mathbb C.
\]
\end{theorem}

\begin{proof}
In the ordered basis $(0,1,2,3)$, use rotation generators
\[
 (J_i)_{jk}=\epsilon_{ijk},\qquad (J_i)_{0\mu}=0,
\]
and boost generators, in the $(\mathbf1,I,J,K)$ convention,
\[
 (K_i)_{0j}=i\delta_{ij},\qquad
 (K_i)_{j0}=-i\delta_{ij},
\]
with all other entries zero.  Solving the linear system
\[
 [C,J_i]=0,\qquad [C,K_i]=0,\qquad i=1,2,3,
\]
for the sixteen complex entries of $C$ gives rank fifteen and hence a
one-complex-dimensional solution space.  The identity is a nonzero solution,
so it spans the commutant.  The exact rational/symbolic rank calculation is
reproduced by
\texttt{verification/su3\_lorentz\_commutant\_check.py}.
\end{proof}

\begin{corollary}[No nontrivial internal $SU(3)$ on the same Lorentz index]
A nontrivial representation of $SU(3)$ acting on the same $W$ index cannot
commute with the full Lorentz action above.  In particular, extending a
Gell--Mann action on
$V=\mathbb C\operatorname{-span}\{I,J,K\}$ by fixing the scalar direction does
not define a Standard-Model internal symmetry on the same carrier: boosts mix
the scalar and spatial directions and fail to commute with generic colour
generators.
\end{corollary}

\section{Meaning for the stabiliser theorem}

The theorem
\[
 \operatorname{Stab}_{GL(V)}(h,\Omega)=SU(3)
\]
remains a correct algebraic statement about the complex three-dimensional
carrier $V$.  The present no-go concerns the next physical identification:
if that same $V$ is literally the spatial part of the Lorentz-transforming
biquaternion $W$, its full $SU(3)$ endomorphism algebra is not an independent
internal symmetry commuting with Lorentz transformations.

Therefore a Standard-Model bridge requires a separation of roles such as
\[
 \mathcal H_{\rm local}
 \sim
 \mathcal H_{\rm Lorentz}\otimes\mathcal H_{\rm internal},
\]
or an equivalent mode/multiplicity construction in which Lorentz acts on the
first factor and $SU(3)$ on the second.  This need not introduce a new
fundamental algebra if the internal multiplicity is derived from existing
$\psi$, spectral, exterior/Fock or other $\Theta$ modes, but the factorisation
must be derived rather than assumed.

\section{Impact on the projector candidate}

A projector
\[
 P=\mathbf1-|n\rangle\langle n|
\]
built directly from the positive Hermitian norm of the same Lorentz carrier is
not automatically Lorentz equivariant.  Indeed the positive
Hilbert--Schmidt/dagger norm is not invariant under nonunitary Lorentz boosts,
as already established in the canonical action audit.

Hence the composite-projector connection is viable as an internal colour
candidate only if its normalized direction $\Phi[\Theta]$ is first shown to be
a Lorentz scalar/internal multiplet with an allowed positive Hermitian pairing.
The naive assignment $\Phi=\Theta$ on the spacetime biquaternion carrier is
not justified.

\begin{tcolorbox}[colback=yellow!8!white,colframe=orange!70!black,title=Status]
\textbf{GAP-SU3-LORENTZ-COMMUTANT: CLOSED AS A CONDITIONAL NO-GO [L1].}
Under the stated same-carrier four-vector representation, no nontrivial internal $SU(3)$ can act on that
$M_2(\mathbb C)$ Lorentz index while commuting with the full Lorentz action.
The UBT action has not yet proved that this is the unique representation carried
by the fundamental $\Theta$ field; representation selection is therefore the
next dynamical subgap.

\medskip
\textbf{Remaining bridge:}
derive a Lorentz-scalar internal triplet/multiplicity bundle from the single
UBT field/mode content and show that the $SU(3)$ stabiliser/operator algebra
acts there.  Only after that separation is established should the local
connection and induced Yang--Mills programme be promoted.
\end{tcolorbox}

\end{document}
